\subsection{I-B08. The logarithmic dense-generator bound already holds in the every-order subclass}
\label{I-B08}

Forbes, Saptharishi, and Shpilka (2014), full version
arXiv:1309.5668, Definition~1.1 (printed p.~3), use the term
\emph{commutative ROABP} for a polynomial admitting a width
bound in \emph{every} variable order.  This is not the class
with one hidden order.  Pairwise commuting layer coefficient
matrices, as in I-26, give a representation in this subclass
by reordering the same factors.

Their Corollary~4.3 (printed p.~22), using Theorem~4.1
(printed pp.~21--22), states that if $f$ has individual degree
less than $d$, admits width-$w$ representations in every order,
and $|F|>nd$, then
\[
 f\ne0\quad\Longleftrightarrow\quad
 f\bigl(\operatorname{SV}_{1+2\lfloor2\log_2w\rfloor}\bigr)
       \ne0.
\]
The stated hitting-set size is
$\operatorname{poly}(n,d)^{O(\log w)}$.
One may use a sufficiently large extension when required.
This known logarithmic pair bound applies to the general
commuting, unipotent representation constructed in I-26.
Consequently its logarithmic necessary pair count matches
the order of a known sufficient bound for that restricted
class.  The new question retained after I-26 concerns the
full one-hidden-order rank-test class, and even a logarithmic
pair bound would still require a separate efficient sampling
argument to give polynomial rather than quasipolynomial size.

The primary definition, theorem, and corollary, including
the local factorization use of the every-order assumption
and the additive composition of two SV maps, were inspected.
The complete article and all of its subsidiary rank-condenser
proofs are not certified here.
